\section{Discussion}
\label{sec:discussion}

\subsection{Key Findings and Their Implications}

Our results support a coherent interpretation running counter to augmentation-invariance theory: \textbf{supervised label correlation is the dominant driver of spurious feature encoding}, not augmentation-based instance-discriminability.

In ERM, the cross-entropy objective trains on a distribution where background texture predicts bird species with 95\% reliability, causing the representation to jointly encode both task and spurious features. MoCo-v3 receives no such label correlation pressure: background IS instance-discriminative (different images have different backgrounds), but without label-correlation amplification, it is encoded less strongly (ratio = 1.027 vs 1.052).

\textbf{ERM $\approx$ DINO} is particularly informative: DINO uses no explicit class labels, yet its ratio (1.050) is statistically indistinguishable from ERM's. This is most parsimoniously explained by DINO's momentum teacher generating class-correlated soft-targets that replicate the label-correlation pressure of supervised ERM at the representation level --- but alternative explanations exist (e.g., DINO's strong augmentations incidentally suppressing non-spurious feature diversity). Testing this via mutual information between DINO teacher-targets and spurious attribute labels is an important future experiment to adjudicate between these explanations.

\textbf{The dataset interaction} reveals that the paradigm $\times$ spurious-attribute relationship is not universal. On CelebA, MoCo-v3 ranks highest --- the opposite of Waterbirds. A plausible interpretation is that gender (distributed across face features, less coarsely instance-discriminative) and background texture (salient, varies dramatically across images) interact differently with each objective's encoding pressure.

\subsection{Limitations}

\textbf{L1: Original directional mechanism refuted.} Our hypothesis predicted contrastive SSL would encode spurious features more strongly. The data show the opposite with large effect size ($d = 5.68$). We treat this as a productive refutation that yields a more nuanced model; the empirical contributions remain valid.

\textbf{L2: h-m1 final statistics incomplete.} The SimCLR-NoBackground mechanism is verified as active ($19\times$ threshold). Full ratio statistics require completing 50-epoch training --- reported as PoC-level finding.

\textbf{L3: No downstream WGA evaluation.} We do not test whether lower spurious ratio translates to better worst-group accuracy via DFR-style fine-tuning. This is well-motivated future work grounded in the DFR framework \citep{Kirichenko2022DFR}.

\textbf{L4: MoCo-v3 as SimCLR proxy.} SimCLR-v2 has no official PyTorch checkpoint; MoCo-v3 was used as the best available contrastive representative. SimCLR-specific validation is future work.

\textbf{L5: Scope.} Results apply to ResNet-50, ImageNet-1k pretraining, Waterbirds and CelebA. ViT backbones, larger pretraining scales, and fine-grained spurious attributes require separate study.

\subsection{Broader Impact}

This work provides backbone selection guidance for fairness-critical applications: contrastive SSL partially suppresses spurious features by design, without requiring group labels or robustness interventions. All paradigms still encode spurious features (ratios $> 1.0$), so the guidance is additive, not a replacement for robustness methods. No new harmful models or datasets are introduced.
